% Complete proofs for P15.21.
\section{Weak normalization with infinite means}
\label{sec:annex-p15-infinite-mean-weak-normalization}

These proofs preserve the positive-summand necessity and truncation
sufficiency of Feller VII.7, Theorem~2
\cite[pp.~236--237]{Feller1971ProbabilityII}, and the distinct
Durrett example calculations
\cite[pp.~65--66]{DurrettPTEAuthor2019}.

\begin{annexitem}{annex:p15.positive-truncation-tools}{thm:p15-positive-truncation-tools}{Tail identities, deletion and the variance estimate}
\begin{proof}
For each fixed nonnegative finite \(x\),
\[
 \min(x,s)=\int_0^s\mathbf1_{\{x>t\}}\,dt,\qquad
 \min(x,s)^2=\int_0^s2t\mathbf1_{\{x>t\}}\,dt.
\]
Tonelli gives the expectations of these nonnegative quantities.
The differences between them and the corresponding deletion
truncations are \(s\mathbf1_{\{x>s\}}\) and
\(s^2\mathbf1_{\{x>s\}}\). Subtracting these finite quantities
proves both identities, including the endpoint atom convention.
The source's integration-by-parts route gives the same identities
directly from Theorem~\ref{thm:p15-feller-integration-by-parts}.
Apply it on \((0,s]\) first to \(u(x)=x\) and then to
\(u(x)=x^2\). The boundary contribution at zero is zero, even if
\(F(0)>0\), and the results are
\[
 \mu(s)=sF(s)-\int_0^s F(t)\,dt,\qquad
 v(s)=s^2F(s)-2\int_0^s tF(t)\,dt.
\]
Substituting \(F(t)=1-q(t)\) gives the stated formulas.
Pointwise, \(x^2\mathbf1_{\{x\leq s\}}\leq
sx\mathbf1_{\{x\leq s\}}\), so \(v(s)\leq s\mu(s)\).

Put \(T_n=\sum_{k=1}^nX_k\mathbf1_{\{X_k\leq s\}}\).
The probability that \(S_n\ne T_n\) is at most \(nq(s)\).
The truncated summands are independent and bounded, so
\(\mathbb ET_n=n\mu(s)\) and
\(\operatorname{Var}(T_n)\leq nv(s)\).
Chebyshev at threshold \(\varepsilon a\), followed by the union
bound, proves the claimed estimate. This is the positive-gain
specialization of Theorem~\ref{thm:p15-feller-truncation-inequality}.

For \(s\to\infty\), the random variables
\((X/s)\mathbf1_{\{X\leq s\}}\) are bounded by one and tend
pointwise to zero. Bounded convergence proves \(\mu(s)/s\to0\).
Finally suppose \(\delta(s)\to0\). For any \(\eta>0\), choose
\(A\geq s_0\) with \(tq(t)\leq\eta\mu(t)\) for \(t\geq A\).
For \(s\geq A\), the second-moment identity and monotonicity of
\(\mu\) give
\[
 0\leq v(s)\leq
 2\int_0^A tq(t)\,dt+2\eta\int_A^s\mu(t)\,dt
 \leq C_A+2\eta s\mu(s).
\]
Here \(C_A<\infty\), since \(q\leq1\), and
\(s\mu(s)\geq s\mu(s_0)\to\infty\).
Divide by \(s\mu(s)\), take the upper limit, and then let
\(\eta\downarrow0\).
\end{proof}
\end{annexitem}

\begin{annexitem}{annex:p15.relative-normalizer-properties}{thm:p15-relative-normalizer-properties}{Growth and comparison without monotonicity}
\begin{proof}
Since \(\{X>0\}=\bigcup_{j\geq1}\{X\geq1/j\}\) has positive
probability, choose \(d>0\) with \(p=\mathbb P(X\geq d)>0\).
Let \(B_n=\sum_{k=1}^n\mathbf1_{\{X_k\geq d\}}\).
Independence gives \(\mathbb EB_n=np\) and
\(\operatorname{Var}(B_n)=np(1-p)\). Chebyshev yields
\(\mathbb P(B_n\geq np/2)\to1\).
As \(S_n\geq dB_n\), the event \(S_n\geq dpn/2\) has
probability tending to one. If \(S_n/a_n\to1\), the event
\(S_n\leq2a_n\) also has probability tending to one.
Their intersection has positive probability eventually, which
forces \(a_n\geq dpn/4\). In particular \(a_n\to\infty\).

The finiteness and common law of \(X_{n+1}\) imply
\(\mathbb P(X_{n+1}>\varepsilon a_n)\to0\) for every
\(\varepsilon>0\). The union bound applied to
\[
 S_{n+1}/a_n-1=(S_n/a_n-1)+X_{n+1}/a_n
\]
proves \(S_{n+1}/a_n\to1\) in probability.
Also \(S_{n+1}/a_{n+1}\to1\) by assumption.
Fix \(0<\eta<1\). Eventually the two events placing these two
ratios in \((1-\eta,1+\eta)\) have an intersection of positive
probability. On that intersection their quotient is the
deterministic number \(a_{n+1}/a_n\), so
\[
 \frac{1-\eta}{1+\eta}<
 \frac{a_{n+1}}{a_n}<
 \frac{1+\eta}{1-\eta}.
\]
Letting \(\eta\downarrow0\) proves the adjacent-ratio limit.
Exactly the same two-event argument with \(S_n/a_n,S_n/b_n\)
proves uniqueness of the relative scale.

For the converse comparison, if \(b_n/a_n\to1\), then \(a_n/b_n\)
is eventually bounded and tends to one. The identity
\[
 S_n/b_n-1=(a_n/b_n)(S_n/a_n-1)+(a_n/b_n-1)
\]
and the definition of convergence in probability give the result.
This entire argument uses the elementary Bernoulli variance
calculation; no strong law is used.
\end{proof}
\end{annexitem}

\begin{annexitem}{annex:p15.positive-normalizer-construction}{thm:p15-positive-normalizer-construction}{The exact first crossing after a positive anchor}
\begin{proof}
Existence of \(s_0\) follows from
\(\{X>0\}=\bigcup_{j\geq1}\{1/j\leq X\leq j\}\).
Some event in that union has positive probability, giving a
positive finite truncated mean. The integer \(N_0\) exists
because \(\mu(s_0)>0\).

Fix \(n\geq N_0\), and write \(h(s)=n\mu(s)/s\) for \(s>0\).
The preceding lemma gives \(h(s)\to0\), so the crossing set is
nonempty. On any bounded interval, bounded convergence for the
truncations shows that \(\mu\) is right-continuous, with left
limit \(\mu(s-)=\mathbb E[X\mathbf1_{\{X<s\}}]\leq\mu(s)\).
Thus \(h\) is right-continuous and its jumps are nonnegative.
Since \(h(s_0)>1\), right continuity excludes crossing points
in some interval \([s_0,s_0+\varepsilon)\); hence \(s_n>s_0\).

Every \(s\in[s_0,s_n)\) has \(h(s)>1\). Taking the left limit
gives \(h(s_n-)\geq1\), whence \(h(s_n)\geq1\).
By the definition of the infimum there are crossing points
in \([s_n,s_n+1/j)\) for every sufficiently large \(j\).
Right continuity then implies \(h(s_n)\leq1\).
Consequently \(h(s_n)=1\), or \(n\mu(s_n)=s_n\).
Monotonicity of \(\mu\) gives \(s_n\geq n\mu(s_0)\).

For \(n+1\), the crossing set is contained in that for \(n\),
since \((n+1)\mu(s)/s\leq1\) implies \(n\mu(s)/s\leq1\).
Their infima therefore satisfy \(s_{n+1}\geq s_n\).
The lower anchor is part of the construction: no values of
\(\mu\) below the positive anchor are used to select a crossing.
\end{proof}
\end{annexitem}

\begin{annexitem}{annex:p15.positive-tail-equivalences}{thm:p15-positive-tail-equivalences}{The complete slow-variation and atom equivalences}
\begin{proof}
Assume \(\delta(s)\to0\). For fixed \(c>1\),
\[
 0\leq\mu(cs)-\mu(s)
   =\mathbb E[X\mathbf1_{\{s<X\leq cs\}}]
   \leq cs\,q(s)=c\delta(s)\mu(s).
\]
Divide by \(\mu(s)>0\). This proves the slow-variation ratio
for \(c>1\); for \(0<c<1\), apply that result with multiplier
\(1/c\) at \(cs\), then take reciprocals. The case \(c=1\)
is immediate.

Conversely suppose \(\mu\) is slowly varying. Fix \(0<\eta<1\).
For all sufficiently large \(u\),
\(\mu(2u)\leq(1+\eta)\mu(u)\).
Consequently, for every integer \(j\geq0\) and sufficiently
large \(s\),
\[
 \mu(2^{j+1}s)-\mu(2^js)
       \leq\eta(1+\eta)^j\mu(s).
\]
The events \(2^js<X\leq2^{j+1}s\), \(j\geq0\), partition
\(\{X>s\}\), because \(X\) is finite. On each such event,
\(X>2^js\); hence
\[
 \begin{aligned}
 s q(s)
 &\leq\sum_{j=0}^\infty
        2^{-j}\{\mu(2^{j+1}s)-\mu(2^js)\}\\
 &\leq\eta\mu(s)\sum_{j=0}^\infty((1+\eta)/2)^j
   =\frac{2\eta}{1-\eta}\mu(s).
 \end{aligned}
\]
Divide by \(\mu(s)\), and then let \(\eta\downarrow0\).
This proves the entire required converse, without a referral to
general regular-variation theory.

The exact atom identities are
\[
 \mu(s)-\mu_<(s)=s\mathbb P(X=s),\qquad
 q_{\geq}(s)=q(s)+\mathbb P(X=s).
\]
Under \(\delta\to0\), for sufficiently large \(s\),
\[
 0\leq\frac{s\mathbb P(X=s)}{\mu(s)}
 \leq\frac{s q(s/2)}{\mu(s)}
 \leq2\delta(s/2)\longrightarrow0.
\]
Thus \(\mu_<(s)/\mu(s)\to1\), making \(\mu_<\) positive
eventually, and
\(s q_{\geq}(s)/\mu_<(s)\to0\).
Conversely, if this last ratio tends to zero, then
\[
 0\leq\delta(s)
 \leq\frac{s q_{\geq}(s)}{\mu_<(s)}
\]
on its eventual positive domain. This proves all equivalences.
The proof does not require \(\mu(s)\to\infty\), so it also
establishes the finite-mean extension of the cited specialization.
\end{proof}
\end{annexitem}

\begin{annexitem}{annex:p15.feller-positive-weak-normalization}{thm:p15-feller-positive-weak-normalization}{Feller's sufficiency and positive-summand necessity}
\begin{proof}
If nonzero real \(a_n\) normalize \(S_n\) to one, an infinite
subsequence with \(a_n<0\) would have \(S_n/a_n\leq0\) almost
surely and therefore error probability one at tolerance \(1/2\).
Hence such a sequence is eventually positive. Altering its finite
initial segment reduces it to the positive domain in the theorem.

Suppose first that \(\rho(s)\to\infty\), or \(\delta(s)\to0\).
Use the constructed \(s_n\) with \(s_n=n\mu(s_n)\).
The two-error estimate at \(s=a=s_n\) gives
\[
 \mathbb P(|S_n/s_n-1|>\varepsilon)
 \leq \frac{n v(s_n)}{\varepsilon^2s_n^2}+nq(s_n)
 =\frac{v(s_n)}{\varepsilon^2s_n\mu(s_n)}+\delta(s_n).
\]
Both terms tend to zero by the truncated-moment lemma and
\(s_n\to\infty\). This supplies the integration-by-parts estimate
and all limits in the source's sufficiency proof.

Conversely, suppose \(S_n/a_n\to1\) in probability for positive
deterministic \(a_n\). Set
\[
 \mu_n=\mu(2a_n),\qquad r_n=\frac{n\mu_n}{a_n}.
\]
The relative-normalizer lemma gives \(a_n\to\infty\) and
\(a_{n+1}/a_n\to1\). Positivity implies
\[
 \mathbb P(S_n<2a_n)
 \leq\mathbb P\!\left(\max_{k\leq n}X_k\leq2a_n\right)
 =[1-q(2a_n)]^n
 \leq e^{-nq(2a_n)}.
\]
The left side tends to one. The inequality
\(1-u\leq e^{-u}\) for \(0\leq u\leq1\), obtained by
differentiating \(e^{-u}-(1-u)\), then forces
\(nq(2a_n)\to0\).

Truncate at \(2a_n\), so \(v(2a_n)\leq2a_n\mu_n\).
The one-sided consequence of the two-error bound is, for every
\(\varepsilon>0\),
\[
 \mathbb P(S_n>n\mu_n+\varepsilon a_n)
       \leq \frac{2r_n}{\varepsilon^2}+nq(2a_n).
\]
If a subsequence had \(r_n\to0\), take \(\varepsilon=1/2\).
Its right side would tend to zero, whereas its left side tends
to one: eventually \(r_n+1/2<3/4\), and
\(\mathbb P(S_n/a_n>3/4)\to1\).
Thus \(\liminf_n r_n>0\). It follows that
\[
 \delta(2a_n)=\frac{2nq(2a_n)}{r_n}\longrightarrow0.
\]
This is the source's necessity argument; it uses the largest
positive summand and the same truncated variance bound as the
source, rather than replacing it by a symmetrization argument.

It remains to prove the limit at every real threshold, without
assuming the chosen \(a_n\) is monotone. Fix a sufficiently large
integer \(N\). For large \(s\), let \(j=j(s)\) be the first
integer \(j\geq N\) for which \(2a_j>s\).
It exists because \(a_j\to\infty\); moreover \(j(s)\to\infty\)
as \(s\to\infty\). Eventually \(j>N\), and the first-crossing
choice gives \(2a_{j-1}\leq s<2a_j\).
The monotonicity of \(\mu\) and the reverse monotonicity of \(q\)
therefore imply
\[
 0\leq\delta(s)
 \leq\frac{2a_j q(2a_{j-1})}{\mu(2a_{j-1})}
 =\frac{a_j}{a_{j-1}}\delta(2a_{j-1})\longrightarrow0.
\]
All denominators are eventually positive. The adjacent-ratio
lemma supplies the remaining factor's limit one. This proves
necessity over the full real parameter domain and completes both
directions. All arguments allow an atom at zero provided
\(\mathbb P(X>0)>0\).
\end{proof}
\end{annexitem}

\begin{annexitem}{annex:p15.positive-normalized-centers}{thm:p15-positive-normalized-centers}{Asymptotic centers at every admissible scale}
\begin{proof}
The relative-normalizer comparison lemma proves the stated
characterization \(a_n/s_n\to1\).
We next justify the truncated-mean comparison at these moving
arguments using the monotonicity that accompanies slow variation.
For any fixed \(0<\eta<1\), eventually
\((1-\eta)s_n\leq a_n\leq(1+\eta)s_n\).
Thus
\[
 \frac{\mu((1-\eta)s_n)}{\mu(s_n)}
 \leq\frac{\mu(a_n)}{\mu(s_n)}
 \leq\frac{\mu((1+\eta)s_n)}{\mu(s_n)}.
\]
Both outer ratios tend to one by the established slow variation.
Consequently
\[
 \frac{n\mu(a_n)}{a_n}
 =\frac{\mu(a_n)}{\mu(s_n)}\frac{s_n}{a_n}\longrightarrow1.
\]
The endpoint equivalence gives the assertion for \(\mu_<\).
Furthermore
\[
 nq(a_n)=\delta(a_n)\frac{n\mu(a_n)}{a_n}\longrightarrow0,
\]
and the proved strict/closed equivalence gives \(nq_{\geq}(a_n)\to0\).

If \(0<m=\mathbb EX<\infty\), then monotone convergence gives
\(\mu(s)\to m\). Also
\[
 s q(s)\leq\mathbb E[X\mathbf1_{\{X>s\}}]\longrightarrow0
\]
by dominated convergence. Thus the criterion holds and
\(s_n/n=\mu(s_n)\to m\).
Comparison of normalizers now shows that \(nm\) works.
If \(\mathbb EX=\infty\), monotone convergence instead gives
\(\mu(s_n)\to\infty\), so \(s_n/n\to\infty\).
These conclusions follow from the weak-normalization construction.
\end{proof}
\end{annexitem}

\begin{annexitem}{annex:p15.positive-normalization-boundaries}{thm:p15-positive-normalization-boundaries}{Computed densities and the independence countermodel}
\begin{proof}
For \(X=c\) and \(s\geq c\), \(\mu(s)=c\).
The crossing set is therefore \([nc,\infty)\) for \(n\geq2\),
giving \(s_n=nc\). For the equally weighted law at \(1,3\),
\(\mu(s)=2\) for \(s\geq3\), so the anchored crossing set
is \([2n,\infty)\) for \(n\geq2\).

For the critical density, direct integration gives
\(\int_1^\infty x^{-2}\,dx=1\),
\(q(s)=\int_s^\infty x^{-2}\,dx=1/s\), and
\[
 \mu(s)=\int_1^s x^{-1}\,dx=\log s,\qquad
 v(s)=\int_1^s1\,dx=s-1.
\]
The truncated first moment diverges as \(s\to\infty\), so the mean is
infinite, whereas \(\delta(s)=1/\log s\to0\).
The crossing construction yields \(s_n=n\log s_n\).
For a direct comparison, take \(b_n=n\log n\), for \(n\geq2\),
and truncate at \(b_n\). Then
\[
 nq(b_n)=1/\log n\to0,\qquad
 \frac{nv(b_n)}{b_n^2}\leq\frac1{\log n}\to0,\qquad
 \frac{n\mu(b_n)}{b_n}
   =1+\frac{\log\log n}{\log n}\to1.
\]
The two-error inequality gives \(S_n/b_n\to1\); comparison of
relative normalizers gives \(s_n\sim b_n\).
For completeness, \(\log u/u\to0\) follows for \(u\geq1\)
from \(\log u=\int_1^u t^{-1}dt\leq2(\sqrt u-1)\);
this justifies the logarithmic ratio used here.

For \(0<\alpha<1\),
\(\int_1^\infty\alpha x^{-\alpha-1}dx=1\),
\(q(s)=s^{-\alpha}\), and
\[
 \mu(s)=\int_1^s\alpha x^{-\alpha}dx
       =\frac{\alpha}{1-\alpha}(s^{1-\alpha}-1).
\]
The moment diverges and division by \(sq(s)=s^{1-\alpha}\)
gives the displayed finite limit of \(\rho\). The necessity
direction excludes every relative weak normalizer.

For the dependence countermodel, let \(X_k=X\) for all \(k\),
where \(X\) takes \(1,3\) with probabilities \(1/2\).
For a proposed normalizer put \(t_n=a_n/n>0\). The event
\(|S_n/a_n-1|\leq1/4\) would require
\[
 t_n\in[4/5,4/3]\quad\hbox{when }X=1,\qquad
 t_n\in[12/5,4]\quad\hbox{when }X=3.
\]
The intervals are disjoint, so its probability is at most \(1/2\)
for every \(n\). Yet the marginal mean is two and its tail
criterion holds. This shows precisely where independence matters.
If \(X=0\) almost surely, every \(S_n/a_n=0\); it cannot
converge in probability to one.
\end{proof}
\end{annexitem}

\begin{annexitem}{annex:p15.cauchy-weak-law-failure}{thm:p15-cauchy-weak-law-failure}{The source tail obstruction and a direct convolution proof}
\begin{proof}
The elementary antiderivative
\(\int (a^2+x^2)^{-1}dx=a^{-1}\arctan(x/a)\) and the
limits \(\arctan(\pm\infty)=\pm\pi/2\) give
\(\int_{\mathbb R}f_a(x)\,dx=1\).
For \(t\geq x>0\),
\[
 \frac{1}{1+x^{-2}}\,t^{-2}
       \leq\frac1{1+t^2}\leq t^{-2}.
\]
Integrating from \(x\) to infinity yields
\[
 \frac{2}{\pi x(1+x^{-2})}
 \leq \mathbb P(|X|>x)\leq\frac2{\pi x}.
\]
The squeeze proves the source tail asymptotic and
\(x\mathbb P(|X|>x)\to2/\pi\).
The weak-tail criterion of
Theorem~\ref{thm:p15-iid-weak-tail-criterion} therefore excludes
every deterministic \(n\)-scale center. Moreover
\(\mathbb E|X|=(2/\pi)\int_0^\infty x/(1+x^2)\,dx=\infty\),
as the antiderivative is \(\tfrac12\log(1+x^2)\).

We also prove the exact law of the average stated in the source.
The established independent-density convolution theorem reduces
the issue to an ordinary integral. Fix \(a,b>0\) and \(x\ne0\);
put
\[
 D=(x^2+(a+b)^2)(x^2+(a-b)^2),\quad
 A=\frac{2x}{D},\quad
 B=\frac{x^2+a^2-b^2}{D},\quad
 C=\frac{x^2-a^2+b^2}{D}.
\]
Here \(D>0\). Multiplication by the two denominators verifies
the rational identity
\[
 \frac1{(y^2+b^2)((y-x)^2+a^2)}
 =\frac{Ay+B}{y^2+b^2}
   +\frac{-A(y-x)+C}{(y-x)^2+a^2}.
\]
Integrate this finite identity over \([-R,R]\), then let
\(R\to\infty\). The first odd numerator has integral zero.
For the shifted odd numerator, the antiderivative is
\(\tfrac12\log((y-x)^2+a^2)\), whose difference at
\(R,-R\) tends to zero. The constant terms tend to
\(\pi B/b\) and \(\pi C/a\). The left integrand is
nonnegative and integrable, so this passage computes its
ordinary full integral; it does not assume separate absolute
integrability of the odd summands.
Multiplying by \(ab/\pi^2\) gives
\[
 (f_a*f_b)(x)=\frac{aB+bC}{\pi}
 =\frac{(a+b)(x^2+(a-b)^2)}{\pi D}
 =\frac{a+b}{\pi(x^2+(a+b)^2)}.
\]
The calculation applies for every \(x\ne0\). Densities may be
changed at the one null point \(x=0\), so it proves equality
of the corresponding probability laws for all \(a,b>0\),
including \(a=b\). No exceptional parameter is discarded.

Induction using independence of disjoint groups now gives
the density \(f_n\) for \(S_n\). The proved affine density
formula gives the density \(n f_n(nx)=f_1(x)\) for \(S_n/n\).
Thus the exact distributional assertion is established directly.
\end{proof}
\end{annexitem}

\begin{annexitem}{annex:p15.petersburg-array-calculation}{thm:p15-petersburg-array-calculation}{The source array with an explicit integer cutoff choice}
\begin{proof}
The masses normalize because \(\sum_{j\geq1}2^{-j}=1\).
In the coin realization the first \(j-1\) tosses must be tails
and the \(j\)-th a head, an event of probability \(2^{-j}\).
The probability of no head in the first \(r\) tosses is \(2^{-r}\);
continuity from above makes the event of no head at all null.
The expectation is
\(\sum_{j\geq1}2^j2^{-j}=\sum_{j\geq1}1=\infty\).

For every integer \(m\geq1\), the geometric sums give
\[
 \mathbb P(X>2^m)=2^{-m},\qquad
 \mathbb P(X\geq2^m)=2^{1-m},
\]
and the truncated moments are exactly
\[
 \mathbb E[X\mathbf1_{\{X\leq2^m\}}]=m,\qquad
 \mathbb E[X^2\mathbf1_{\{X\leq2^m\}}]
       =\sum_{j=1}^m2^j=2^{m+1}-2.
\]
The atom at \(2^m\) accounts for the difference between the
two tails. With \(m=m(n)\), multiplication of the tails by \(n\)
gives the displayed \(2^{-K(n)}\) and \(2^{1-K(n)}\).
For the row \(X_{n,k}=X_k\), \(1\leq k\leq n\),
independence is inherited from the original sequence.
Furthermore
\[
 b_n\to\infty,\quad
 \sum_{k=1}^n\mathbb P(X_{n,k}>b_n)=2^{-K(n)}\to0,\quad
 b_n^{-2}\sum_{k=1}^n\mathbb E\overline X_{n,k}^{\,2}
          \leq 2n/b_n=2^{1-K(n)}\to0.
\]
These are all hypotheses of
Theorem~\ref{thm:p15-triangular-array-weak-law}; its center
is exactly \(a_n=nm(n)\). Explicitly it gives, for \(t>0\),
\[
 \mathbb P(|S_n-a_n|>t b_n)
        \leq2^{-K(n)}+t^{-2}2^{1-K(n)}\longrightarrow0.
\]
Using the larger non-strict tail, as the source does, is also
valid because it bounds the deletion probability from above.

For any such choice with \(K(n)/\log_2n\to0\), the exact identity
\(a_n/(n\log_2n)=1+K(n)/\log_2n\) proves the separate
deterministic center limit asserted in the source.
If \(K(n)\leq\log_2\log_2n\) eventually, then
\[
 0<b_n/(n\log_2n)=2^{K(n)}/\log_2n\leq1.
\]
Therefore the centered conclusion remains true with denominator
\(n\log_2n\). Also \(K(n)\to\infty\) makes \(K(n)>0\)
eventually, and
\[
 0\leq K(n)/\log_2n
       \leq(\log_2\log_2n)/\log_2n\longrightarrow0.
\]
The last limit follows from \(\log u/u\to0\), proved above,
with a constant change of logarithmic base. Consequently
\(a_n/(n\log_2n)=1+K(n)/\log_2n\to1\).
Adding this deterministic center difference proves the limit.

For the explicit choice, the floor inequality gives
\[
 \log_2\log_2n-1<K(n)\leq\log_2\log_2n
                                      \qquad(n\geq2).
\]
Thus \(K(n)\to\infty\), the required upper bound holds, and
\(m(n)\) is an integer by construction. This supplies an actual
simultaneous choice of every parameter used in the source route.
\end{proof}
\end{annexitem}

\begin{annexitem}{annex:p15.petersburg-normalizers-and-prices}{thm:p15-petersburg-normalizers-and-prices}{The separate criterion calculation and numerical scales}
\begin{proof}
For \(s\geq2\), the retained positive masses are exactly those
with \(1\leq j\leq\ell(s)\). The preceding finite sums therefore
give the displayed \(\mu(s)\) and \(v(s)\); summing the remaining
geometric tail gives \(q(s)=2^{-\ell(s)}\).
Since \(2^{\ell(s)}\leq s<2^{\ell(s)+1}\),
\[
 \delta(s)=\frac{s\,2^{-\ell(s)}}{\ell(s)}
             <\frac2{\ell(s)}\longrightarrow0.
\]
Feller's theorem applies, and its exact equation is
\(s_n=n\ell(s_n)\). The explicit triangular-array proof has
already supplied \(n\log_2n\) as another relative weak
normalizer. Their ratio tends to one by uniqueness of scale.
This preserves the two source routes rather than using the
criterion as a substitute for the source array calculation.

At \(n=1024=2^{10}\), \(n\log_2n=10240\), whose base-two
logarithm lies between \(13\) and \(14\).
Hence \(m(n)=13\), \(K(n)=3\), \(b_n=2^{13}=8192\), and
\(a_n=1024\cdot13=13312\).
The strict tail is \(1024/8192=1/8\), and
\[
 \frac{1024(2^{14}-2)}{(2^{13})^2}
   =\frac14-\frac1{32768}.
\]
These finite bounds describe this truncation and are not
assertions of a small error at a prescribed tolerance.

The leading per-play scale is
\((n\log_2n)/n=\log_2n\), giving \(10\) at \(n=1024\).
Solving \(\log_2n=40\) gives
\(n=2^{40}=1\,099\,511\,627\,776\).
For every fixed \(c>0\), eventually
\((1/2)\log_2n>c\). On the event
\(S_n/(n\log_2n)>1/2\) one then has \(S_n>cn\), and the
former event has probability tending to one. This proves the
fixed-fee probability statement without assigning an expected
finite price to a gain of infinite expectation.

Finally \(q(2^j)=2^{-j}\), so \(2^jq(2^j)=1\).
The common positive law therefore fails the \(n\)-scale tail
criterion in Theorem~\ref{thm:p15-iid-weak-tail-criterion}.
Its necessary direction excludes all deterministic centers
for \(S_n/n\), even though the larger denominator just proved
gives relative convergence to one.
\end{proof}
\end{annexitem}
